% 分野別類題: 2階線形偏微分方程式の分類と標準形
% A u_xx + 2B u_xy + C u_yy + ... = 0 の型判定・特性曲線を用いた標準形化。
% コンパイル: TEXINPUTS="../../../00_common:$TEXINPUTS" latexmk -xelatex problems.tex
\documentclass[a4paper,11pt]{article}
\usepackage{preamble}

\begin{document}

\kamokumei{類題演習 --- 2階線形偏微分方程式の分類と標準形}

\shijibun{次の2階線形偏微分方程式について、判別式 $B^{2}-AC$ による型（楕円型・放物型・双曲型）を判定し、必要に応じて特性曲線を用いて標準形へ変換し一般解を求めよ。ここで方程式は
\[
A\,u_{xx} + 2B\,u_{xy} + C\,u_{yy} + (\text{低階の項}) = 0
\]
の形で与えられ、$A,B,C$ は $x,y$ の関数（定数を含む）である。}

\shomon{問1.}{次の方程式の型を判定せよ（領域によって型が変わる場合はその境界も述べよ）。}
\[
y\,u_{xx} + u_{yy} = 0
\]

\shomon{問2.}{次の双曲型方程式を特性曲線を用いて標準形に変換し、一般解を求めよ。}
\[
u_{xx} - 4u_{xy} + 3u_{yy} = 0
\]

\shomon{問3.}{次の放物型方程式を標準形に変換し、初期条件 $u(x,0) = x^{2}$, $u_{y}(x,0) = 0$ を満たす解を求めよ。}
\[
u_{xx} - 2u_{xy} + u_{yy} = 0
\]

\end{document}
